% ---------------------------------------------------------------------------------------------------- 11
\section{Conclusions}\label{sec:conclusions}
\begin{itemize}
  \item \textbf{Early exit.} The fastest rule changes with batch size and input length; a rule derived from the
  measured cost tracks it. At matched accuracy it leads on QQP (2--17\% on the GTX 1650 cost table; 8--14\% timed on
  the A6000, replaying recorded decisions on a random-weight network) and on Llama classifiers with trained weights
  (3.66\x). Certified selection keeps the one-point budget on re-splits. The slowdown bound is proved for the trimmed
  collective stop and held empirically for the full rule in every tested batch. No gain on RoBERTa MNLI or on LLM
  decoders, and the accuracy loss lands on hard inputs.
  \item \textbf{Minrank.} $\operatorname{rank}_{\Fp}(I+\varepsilon D)=\binom{n-1}{k-1}$ exactly when $q_p(k)\mid n-2k$.
  Hence, whenever $q_p(k)\mid n-2k$, $\operatorname{minrk}_{\Fp}(K(n,k))=\beta=\alpha$ with an explicit optimal
  index code; this is new for $k\nmid n$. The converse does not hold for the minrank: $K(9,3)$ over $\mathbb{F}_2$ has
  $\operatorname{minrk}=\alpha$ by Baranyai's theorem although $4\nmid3$. Small cases such as $K(9,2)$ over
  $\mathbb{F}_2$ remain open.
  \item \textbf{Agent.} A Solver--Critic--Verifier pipeline, not given the graph's name, proves $\beta$ on 19 of 25
  graphs with each of the four Claude settings run on the full benchmark; the no-LLM Kneser script also proves 19. The
  LLM's measurable contribution is recognising the Kneser family, which a no-LLM fast path also does (17 of the 23
  graphs it ran on, all that the Claude agents prove there). The pipeline's speed on large graphs (43--55\,s against
  88\,s at 10{,}626 receivers) comes from its $O(N^2r)$ rank engine; the LLM adds a median of 30\,s per graph (12--50\,s). Whether an LLM
  helps on graphs that no fixed rule recognises is untested.
\end{itemize}

% ---------------------------------------------------------------------------------------------------- references
\begin{thebibliography}{99}\small
\bibitem{deebert} J.~Xin, R.~Tang, J.~Lee, Y.~Yu, J.~Lin. DeeBERT: Dynamic Early Exiting for Accelerating BERT Inference. ACL 2020.
\bibitem{pabee} W.~Zhou, C.~Xu, T.~Ge, J.~McAuley, K.~Xu, F.~Wei. BERT Loses Patience: Fast and Robust Inference with Early Exit. NeurIPS 2020.
\bibitem{drex} X.~Liu et al. Dynamic Rebatching for Efficient Early-Exit Inference with DREX. arXiv:2512.15705, 2025.
\bibitem{geep} L.~Li et al. Accelerating BERT Inference with GPU-Efficient Exit Prediction. Frontiers of Computer Science 18(3), 2024.
\bibitem{miao} R.~Miao, Y.~Yan, X.~Yao, T.~Yang. An Efficient Inference Framework for Early-exit Large Language Models. arXiv:2407.20272, 2024.
\bibitem{apparate} Y.~Dai, R.~Pan, A.~Iyer, K.~Li, R.~Netravali. Apparate: Rethinking Early Exits to Tame Latency-Throughput Tensions in ML Serving. SOSP 2024.
\bibitem{fastsafe} M.~Jazbec, A.~Timans, et al. Fast yet Safe: Early-Exiting with Risk Control. NeurIPS 2024.
\bibitem{e3} A.~Iyer et al. Improving DNN Inference Throughput Using Practical, Per-Input Compute Adaptation (E3). SOSP 2024.
\bibitem{sloth} S.~Hong, Y.~Kaya, I.-V.~Modoranu, T.~Dumitra\c{s}. A Panda? No, It's a Sloth: Slowdown Attacks on Adaptive Multi-Exit Neural Network Inference. ICLR 2021.
\bibitem{lostpatience} Z.~Coalson, G.~Ritter, R.~Bobba, S.~Hong. BERT Lost Patience Won't Be Robust to Adversarial Slowdown. NeurIPS 2023.
\bibitem{falsesense} Y.~Chen et al. Dynamic Transformers Provide a False Sense of Efficiency. ACL 2023.
\bibitem{chow} Y.~S.~Chow, H.~Robbins, D.~Siegmund. \emph{Great Expectations: The Theory of Optimal Stopping}. Houghton Mifflin, 1971.
\bibitem{ltt} A.~N.~Angelopoulos, S.~Bates, E.~J.~Cand\`es, M.~I.~Jordan, L.~Lei. Learn then Test: Calibrating Predictive Algorithms to Achieve Risk Control. arXiv:2110.01052, 2021.
\bibitem{labeleff} N.~Cesa-Bianchi, G.~Lugosi, G.~Stoltz. Minimizing Regret with Label Efficient Prediction. IEEE Trans.\ Information Theory 51(6), 2005.
\bibitem{ceebert} D.~J.~Bajpai, M.~K.~Hanawal. CeeBERT: Cross-Domain Inference in Early Exit BERT. Findings of ACL 2024.
\bibitem{ucbee} R.~G.~Pacheco et al. UCBEE: A Multi-Armed Bandit Approach for Early-Exit in Neural Networks. IEEE Trans.\ Network and Service Management, 2024.
\bibitem{roberta} Y.~Liu et al. RoBERTa: A Robustly Optimized BERT Pretraining Approach. arXiv:1907.11692, 2019.
\bibitem{mnli} A.~Williams, N.~Nangia, S.~Bowman. A Broad-Coverage Challenge Corpus for Sentence Understanding through Inference. NAACL 2018.
\bibitem{glue} A.~Wang et al. GLUE: A Multi-Task Benchmark and Analysis Platform for Natural Language Understanding. ICLR 2019.
\bibitem{vit} A.~Dosovitskiy et al. An Image Is Worth 16x16 Words: Transformers for Image Recognition at Scale. ICLR 2021.
\bibitem{cifar} A.~Krizhevsky. Learning Multiple Layers of Features from Tiny Images. Tech.\ report, University of Toronto, 2009.
\bibitem{qwen} Qwen Team. Qwen2.5 Technical Report. arXiv:2412.15115, 2024.
\bibitem{llama} Llama Team, Meta AI. The Llama 3 Herd of Models. arXiv:2407.21783, 2024.
\bibitem{wikitext} S.~Merity, C.~Xiong, J.~Bradbury, R.~Socher. Pointer Sentinel Mixture Models. ICLR 2017.
\bibitem{skipdecode} L.~Del Corro et al. SkipDecode: Autoregressive Skip Decoding with Batching and Caching for Efficient LLM Inference. arXiv:2307.02628, 2023.
\bibitem{layerskip} M.~Elhoushi et al. LayerSkip: Enabling Early Exit Inference and Self-Speculative Decoding. ACL 2024.
\bibitem{calm} T.~Schuster et al. Confident Adaptive Language Modeling. NeurIPS 2022.
\bibitem{baryossef} Z.~Bar-Yossef, Y.~Birk, T.~S.~Jayram, T.~Kol. Index Coding with Side Information. IEEE Trans.\ Information Theory 57(3), 2011.
\bibitem{gh} A.~Golovnev, I.~Haviv. The (Generalized) Orthogonality Dimension of (Generalized) Kneser Graphs: Bounds and Applications. Theory of Computing, 2022.
\bibitem{ducey} J.~E.~Ducey, L.~Engelthaler, J.~Gathje, B.~Jones, I.~Pfaff, J.~Plute. Integer Diagonal Forms for Subset Intersection Relations. arXiv:2310.09227, 2023.
\bibitem{rcps} S.~Bates, A.~Angelopoulos, L.~Lei, J.~Malik, M.~I.~Jordan. Distribution-Free, Risk-Controlling Prediction Sets. Journal of the ACM 68(6), 2021 (Hoeffding--Bentkus bound).
\bibitem{wsr} I.~Waudby-Smith, A.~Ramdas. Estimating Means of Bounded Random Variables by Betting. Journal of the Royal Statistical Society Series B 86(1), 2024.
\bibitem{lora} E.~J.~Hu, Y.~Shen, P.~Wallis, Z.~Allen-Zhu, Y.~Li, S.~Wang, L.~Wang, W.~Chen. LoRA: Low-Rank Adaptation of Large Language Models. ICLR 2022.
\bibitem{funsearch} B.~Romera-Paredes et al. Mathematical Discoveries from Program Search with Large Language Models. Nature 625, 2024.
\bibitem{baranyai} Zs.~Baranyai. On the Factorization of the Complete Uniform Hypergraph. In \emph{Infinite and Finite Sets}, Colloq.\ Math.\ Soc.\ J\'anos Bolyai 10, North-Holland, 1975.
\bibitem{ekr} P.~Erd\H{o}s, C.~Ko, R.~Rado. Intersection Theorems for Systems of Finite Sets. Quarterly Journal of Mathematics 12(1), 1961.
\bibitem{lucas} E.~Lucas. Th\'eorie des fonctions num\'eriques simplement p\'eriodiques. American Journal of Mathematics 1, 1878.
\bibitem{cpsat} L.~Perron, F.~Didier. CP-SAT. Google OR-Tools, \url{https://developers.google.com/optimization}.
\bibitem{cadical} A.~Biere, T.~Faller, K.~Fazekas, M.~Fleury, N.~Froleyks, F.~Pollitt. CaDiCaL 2.0. CAV 2024.
\bibitem{drattrim} N.~Wetzler, M.~J.~H.~Heule, W.~A.~Hunt. DRAT-trim: Efficient Checking and Trimming Using Expressive Clausal Proofs. SAT 2014.
\bibitem{flint} The FLINT team. FLINT: Fast Library for Number Theory. \url{https://flintlib.org}.
\bibitem{tabucol} A.~Hertz, D.~de~Werra. Using Tabu Search Techniques for Graph Coloring. Computing 39(4), 1987.
\bibitem{vllm} W.~Kwon et al. Efficient Memory Management for Large Language Model Serving with PagedAttention. SOSP 2023.
\bibitem{llamacpp} G.~Gerganov et al. llama.cpp. \url{https://github.com/ggml-org/llama.cpp}.
\end{thebibliography}

% ==================================================================================================== appendices
\appendix
\section{Environment, versions and seeds}\label{app:env}

\begin{table}[ht]
  \centering\small
  \begin{tabularx}{\textwidth}{@{}>{\raggedright\arraybackslash}p{3.6cm}X@{}}
    \toprule
    Item & Value \\
    \midrule
    Repository & \repo, branch \texttt{main}, commit \texttt{50f1ef4} (the progress report of 13 November is commit \texttt{14e5bfa}) \\
    GPU server & Linux 6.8.0-151-generic; 2\x{} Intel Xeon Silver 4216 (64 threads); 4\x{} RTX A6000 (48\,GB) and 2\x{} RTX 3080 Ti (every server number in this report is from an RTX A6000); NVIDIA driver 580.159.04, CUDA 13.0 \\
    Laptop (Tables~\ref{tab:cost}--\ref{tab:summary}) & GTX 1650 4\,GB, fp32, 4 CPU threads; package versions not recorded \\
    Agent container (Table~\ref{tab:ablation}) & 4-core container; each pipeline pinned to one core (\texttt{taskset}, \texttt{OMP\_NUM\_THREADS=1}) \\
    Python & 3.12.3 in a virtual environment \\
    Key packages & torch 2.14.0+cu130, transformers 5.17.0, datasets 5.0.1, accelerate 1.15.0, numpy 2.5.3, scipy 1.18.1, pandas 3.0.6, scikit-learn 1.9.1, matplotlib 3.11.2, ortools 9.15.6755, python-sat 1.9.dev24, networkx 3.7, cvxpy 1.9.3, highspy 1.16.0, python-flint, numba \\
    Exact versions & \file{requirements-lock.txt} (\texttt{pip freeze} of the environment used for every run) \\
    External binaries & CaDiCaL and drat-trim (DRAT checks); vLLM and llama.cpp b12780 (local LLMs) \\
    \bottomrule
  \end{tabularx}
  \caption{Hardware and software.}
\end{table}

\textbf{Seeds.}
\begin{itemize}
  \item Training: seed 0 for every model (data order and initialisation of new heads). Tuning/held-out split of each
  trace: a permutation with seed 0.
  \item Selection: fit/cert split with seed 0; disjoint batches with seeds $B$ (fit), $B+1$ (cert), $B+2$ (old
  protocol), $100+B$ (held-out). Random-weight timing network: \texttt{torch.manual\_seed(0)}.
  \item Re-splits: seeds 0--29 (permutation seed $1000+s$, test batches $s+7$). Bandit: seeds 0--9. Load traces: 3
  arrival seeds. Random-search candidates: seed 0.
  \item Agent: the seed sets the vertex relabelling (seed 0 for Table~\ref{tab:ablation}; seeds 0 and 1 for the blind
  ablation).
  \item CPU reproductions were run with \texttt{PYTHONHASHSEED=0} and no visible GPU.
\end{itemize}

\textbf{Fine-tuned checkpoints} are not in the repository (they are large); Step~R\ref{st:train} retrains them. The
validation traces and all measured cost tables are committed, so every CPU step runs without checkpoints.

% ----------------------------------------------------------------------------------------------------
\section{Known issues and corrections}\label{app:issues}
\begin{itemize}
  \item \textbf{Calibration-cache bug (fixed 26 November).} The cost rule's calibration table was cached under the
  memory address of the trace object it was fitted on. In the re-split experiments new traces are built for every
  job; after garbage collection a later job could get the same address and reuse a table fitted on another re-split.
  Only the cost-rule rows of Tables~\ref{tab:resplit} and~\ref{tab:cert} were affected, and the result depended on
  the number of worker processes. The fix keeps a weak reference and refits when the object differs; runs with 3 and
  6 workers now give identical files. No claim changed status. This report shows the fixed values, with the 13
  November values of Table~\ref{tab:resplit} in brackets.
  \item \textbf{Labelling of Tables~\ref{tab:t11}--\ref{tab:t12b}.} These timings replay the trained model's recorded
  decisions on a random-weight fp32 network with random token ids, as on the GTX 1650. Accuracy is exact and timing
  does not depend on weights, but they are not ``trained checkpoint, real inputs'' timings as first labelled.
  Tables~\ref{tab:t15}--\ref{tab:t16} do use trained weights and real prompts.
  \item \textbf{Ten numbers in the 13 November text do not match the committed result files} (two in
  Table~\ref{tab:corrections}, the rest in the running text; none in a table of measurements, and none changes a
  claim). For one of them the file changed after 13 November: the CIFAR-100 robust-rule speedup (1.12\x{} in the
  text) is now a 22 November rerun giving 1.24\x, so the 13 November number may have matched the file of that day.
  Another, the simulator error of 15--32\%, cannot be traced to any file. All are corrected in Part~I and marked with
  footnotes: the simulator error of Table~\ref{tab:corrections} (now the per-policy 1--47\% of the files) and its
  1.12\x{} (file: 1.24\x); the heads' share under CUDA graphs (19\% on MNLI only, 12\%
  on QQP); the single-run bandit slopes (superseded by ten seeds); consensus collapsing at $k=1$ (only on CIFAR-100);
  the frozen Llama heads (8.2 points short, not about 5); and the 5.6 of the LLM ``Expected'' paragraph (4.0).
  \item \textbf{Not run:} the certified CUDA-graph tables (Section~\ref{sec:server-cert}), ViT at $B=16,64,128$, and
  trained checkpoints on the GTX 1650.
\end{itemize}

% ----------------------------------------------------------------------------------------------------
\section{The agents' system prompts (verbatim, prompt version v4)}\label{app:prompts}
The LLM chooses and orders tools, argues and explains; it never writes a bound.

\subsection*{Solver}
\lstinputlisting{prompts/solver.txt}
\subsection*{Critic and checkpoint}
\lstinputlisting{prompts/critic.txt}
\subsection*{Verifier}
\lstinputlisting{prompts/verifier.txt}
\subsection*{Tool list appended to the Solver and Critic prompts}
\lstinputlisting{prompts/tools.txt}

% ----------------------------------------------------------------------------------------------------
\section{Implementation map: where each step is in the repository}\label{app:map}

All paths are relative to the root of \repo{} (branch \texttt{main}, commit \texttt{50f1ef4}), except the sources
of this report, which are not in that commit (last row). To get the code and the environment:
\begin{lstlisting}[language=bash]
git clone https://github.com/yneha3229/Optimal-Batch-Level-Early-Exit.git
cd Optimal-Batch-Level-Early-Exit && git checkout 50f1ef4
python3.12 -m venv venv && venv/bin/pip install -r requirements-lock.txt
\end{lstlisting} Every result file named in the captions of this report is committed
under \file{results/} (Part~I) or \file{kneser-agent/results/} and \file{minrank/results/} (Part~II), with its raw
log under \file{logs/}.

{\small
\renewcommand{\arraystretch}{1.15}
\begin{longtable}{@{}>{\raggedright\arraybackslash}p{3.7cm}>{\raggedright\arraybackslash}p{11.9cm}@{}}
  \toprule
  Step & Implementation (files) \\
  \midrule
  \endhead
  \multicolumn{2}{@{}l}{\textbf{Part I: batch-level early exit}} \\
  R\ref{st:train} Fine-tune with exits & \file{traces/train\_exits.py} (training loop, loss, optimiser);
    \file{simulator/exit\_model.py} (exit heads, model wrapper); \file{configs/train\_\{mnli,qqp,cifar100\}.yaml}
    (hyperparameters); wrappers \file{scripts/train\_text.sh}, \file{scripts/train\_vision.sh} \\
  R\ref{st:traces} Record traces & \file{traces/train\_exits.py} (\texttt{log\_traces}, option \texttt{--traces-only});
    output \file{traces/<task>\_val.parquet} \\
  R\ref{st:profile} Cost profile & \file{experiments/local\_hw.py} subcommand \texttt{profile} (classes \texttt{Device},
    \texttt{TCost}); output \file{results/local\_cost\_<gpu>\_T<len>.json}. Server runs of Figure~\ref{fig:server}:
    \file{runner/gpu\_runner.py} (eager), \file{runner/graph\_runner.py} (CUDA graphs), plot
    \file{analysis/report\_figures.py} \\
  R\ref{st:calib} Calibration tables & \file{simulator/policies.py} (class \texttt{Calibration}) \\
  R\ref{st:policies} Policies and grids & \file{experiments/local\_hw.py} (\texttt{theory\_c}, \texttt{make},
    \texttt{grid}); \file{simulator/policies.py} (\texttt{threshold}, \texttt{majority},
    \texttt{param\_collective}, \texttt{theory\_rule}) \\
  R\ref{st:replay} Replay & \file{experiments/local\_hw.py} (\texttt{replay}); \file{simulator/simulator.py} \\
  R\ref{st:select} Certified selection & \file{experiments/local\_hw.py} subcommand \texttt{select}
    (\texttt{select\_family}, \texttt{pareto\_path}, \texttt{certified}); output \file{results/local\_select\_*.json} \\
  R\ref{st:time} Timed runs & \file{experiments/local\_hw.py} subcommand \texttt{time} (\texttt{Device},
    \texttt{LLMDevice}, \texttt{GraphDevice}); drivers \file{scripts/run\_local.sh} (GTX 1650),
    \file{scripts/run\_server\_a6000.sh}, \file{scripts/run\_server\_a6000\_vit.sh} (A6000); output
    \file{results/local\_time\_*.csv} \\
  R\ref{st:front} Fronts, hard quarter & \file{experiments/local\_hw.py} subcommand \texttt{front};
    \file{scripts/run\_front.sh}; output \file{results/local\_front\_*.csv} \\
  R\ref{st:resplit} Re-splits & \file{experiments/local\_hw.py} subcommand \texttt{resplit}; output
    \file{results/local\_resplit\_*.csv}; cache-fix rerun \file{scripts/rerun\_cache\_fix.sh} \\
  R\ref{st:certcheck} Certificates & \file{experiments/cert\_check.py}, \file{experiments/delta\_check.py};
    \file{scripts/run\_cert\_check.sh}, \file{scripts/run\_mnli\_delta.sh}; output \file{results/cert\_check\_*.csv} \\
  R\ref{st:load} Fluctuating load & \file{experiments/load\_trace.py}; \file{scripts/run\_load\_trace.sh}; output
    \file{results/load\_trace\_*.csv} \\
  R\ref{st:bandit} Bandit & \file{bandit/bandit.py} (\texttt{LabelFreeStopRule}, \texttt{CeeBertStyle}),
    \file{bandit/run\_bandit.py}; \file{scripts/run\_bandit\_seeds.sh}; output \file{results/bandit\_slopes\_*\_10seeds.csv} \\
  R\ref{st:robust} Robustness & \file{attacks/deepsloth\_vit.py} (attack inputs, output
    \file{traces/cifar100\_attack.parquet}); \file{experiments/robust\_bound.py}; \file{simulator/attacks.py};
    \file{scripts/run\_robust\_cert.sh}; tests in \file{tests/} and \file{simulator/test\_simulator.py} \\
  R\ref{st:llm} LLM tracks & \file{decoder/train\_cls\_exit\_lora.py} (classifiers), \file{decoder/train\_exit\_lora.py}
    and \file{decoder/train\_decoder\_exits.py} (decoder exits), \file{decoder/layerskip\_exits.py};
    \file{analysis/adaptivity\_vs\_batch.py}; timing via \file{experiments/local\_hw.py} (\texttt{LLMDevice}) and
    \file{scripts/run\_llm\_cls\_hw.sh} \\
  R\ref{st:theory} Theory checks & \file{theory/verify\_58.py}, \file{experiments/theory58\_front.py}; proofs
    \file{theory/THEORY\_58.md} \\
  Checking a reproduction & \file{analysis/check\_report\_numbers.py} (every printed number against its file);
    \file{scripts/reproduce\_cpu.sh} (regenerate all CPU results and diff, with \file{analysis/diff\_repro.py});
    \file{scripts/reproduce\_gpu.sh} (the GPU steps in order); per-number map \file{docs/REPRODUCE\_EARLY\_EXIT.md};
    the 13 November report generator \file{analysis/make\_report.py} \\
  \midrule
  \multicolumn{2}{@{}l}{\textbf{Part II: Kneser graphs}} \\
  R\ref{st:rank} Theorem by exact ranks & \file{minrank/analysis/verify\_star\_basis.py} \\
  R\ref{st:indep} Independent rank code & \file{minrank/analysis/verify\_independent.py} \\
  R\ref{st:secondproof} Second proof & \file{minrank/analysis/verify\_all\_primes.py},
    \file{minrank/analysis/verify\_general\_k.py}, \file{minrank/analysis/verify\_kneser\_proof.py}; statement and
    proofs in \file{proof.md} \\
  R\ref{st:codecheck} Code check & \file{minrank/analysis/verify\_star\_basis.py}; decoding simulation of the agent's
    verifier \file{kneser-agent/kagent/verify.py} \\
  R\ref{st:sat} SAT and DRAT & \file{minrank/exact/sat.py} ($\mathbb{F}_2$ encoding), \file{minrank/exact/symmetry.py}
    (symmetry breaking), \file{minrank/exact/sat\_fq.py} ($\mathbb{F}_q$), \file{minrank/exact/families.py},
    \file{minrank/exact/field\_sep.py}; DRAT check \file{scripts/minrank/10\_drat\_check.sh}; output
    \file{minrank/results/drat\_checks.csv} \\
  R\ref{st:graphs} Benchmark graphs & \file{kneser-agent/kagent/graphs.py}; lists \file{kneser-agent/bench/specs.json},
    ground truth \file{kneser-agent/bench/truth.json} \\
  R\ref{st:baselines} Baselines & \file{kneser-agent/bench/baselines.py} (standard methods),
    \file{kneser-agent/bench/kushagra\_kneser\_agent.py} (Kneser agent, no LLM); container run
    \file{kneser-agent/bench/ablation/ctr\_baselines.sh} \\
  R\ref{st:agentrun} Agent runs & \file{kneser-agent/kagent/pipeline.py} (orchestrator, scheduler, bound board,
    trace); \file{kagent/agents.py} (prompts, JSON schemas, reply validation, default plan); \file{kagent/llm.py}
    (backends: none, local, file-based); \file{kagent/tools.py} (summary, greedy bounds, tool library);
    \file{kagent/rank.py} (rank engine); \file{kagent/verify.py} (independent verifier); \file{kagent/view.py}
    (terminal view of a trace); \file{kneser-agent/bench/run\_agent.py}; blind runs
    \file{kneser-agent/bench/ablation/full\_run.sh} and \file{next\_req.py} \\
  R\ref{st:score} Scoring and charts & \file{kneser-agent/bench/report.py} (\file{RESULTS.md}),
    \file{bench/ablation/score.py}, \file{bench/ablation/plot\_time.py}, \file{bench/ablation/plot\_deck.py};
    results \file{kneser-agent/results/runs.jsonl}, \file{baselines\_ctr\{,\_xl\}.jsonl} \\
  R\ref{st:blind} Blind ablation & \file{kneser-agent/bench/ablation/start\_cells.sh}, \file{cpu\_grid.sh} (Qwen on
    CPU), \file{score.py}, \file{labels.py} (objection labels, \file{label\_detail.json}), \file{plot.py} \\
  R\ref{st:tests} Unit tests & \file{kneser-agent/tests/} \\
  \midrule
  \multicolumn{2}{@{}l}{\textbf{This report}} \\
  Report & Not part of commit \texttt{50f1ef4}; distributed with this report: \file{main.tex} and its parts, tables
    generated by \file{make\_tables.py}, three figures redrawn from the committed data by \file{make\_figs.py}; build
    with \texttt{python make\_tables.py \&\& python make\_figs.py \&\& latexmk -pdf main.tex} \\
  \bottomrule
\end{longtable}
}
